Universal access to basic sanitation services is essencial to public health and to reducing inequalities in living conditions. In this regard, services such as water supply, sewage collection and treatment, and solid waste collection and proper disposal contribute to a healthy environment.

In Brazil, Law No. 14,026/2020 updated the legal framework for basic sanitation and established that service provision contracts must include targets of providing 99\% of the population with drinking water and 90\% with sewage collection and treatment by December 31, 2033. Given these targets and the diversity of conditions across Brazilian municipalities, monitoring sanitation conditions is fundamental to inform the debate on universal access and identify the municipalities most in need of public policies. Because sanitation encompasses different dimensions, such a comparison requires decisions about which indicators to consider and how to combine them into a score for each municipality.

Among the tools available for this purpose, the 2026 edition of the ABES Ranking of Universal Access to Sanitation compares Brazilian municipalities using official 2024 data from the National Basic Sanitation Information System (SINISA). The ranking considers five indicators:

\begin{enumerate}
    \item \textbf{Population served by a water supply network (IAG0001):} the percentage of the total population served by a water supply network, calculated as the ratio of the population served to the municipality’s total population.
    \item \textbf{Population served by a sewage collection network (IES0001):} the percentage of the total population served by a sewage collection network, calculated as the ratio of the population served to the municipality’s total population.
    \item \textbf{Treated sewage relative to water consumed (IES2003):} the sum of the volumes of treated sewage and raw sewage exported for treatment, expressed as a percentage of the volume of water consumed. In this study, values above 100\% are capped at that same number.
    \item \textbf{Household solid waste collection coverage (IRS0001):} the percentage of the total population served by household solid waste collection, calculated as the ratio of the population served to the municipality’s total population.
    \item \textbf{Proper final disposal of municipal solid waste (adapted IRS3002):} calculated by subtracting from 100\% the percentage of improper final disposal reported by SINISA indicator IRS3002. The indicator is thus oriented so that higher values represent better disposal conditions.
\end{enumerate}

Based on these indicators, ABES produces its ranking by adding the five values, giving each indicator the same nominal contribution to a total score ranging from 0 to 500 points. This approach involves choices about the importance assigned to each dimension and how performance across dimensions is combined. For example, when the indicators are added together, higher performance in one dimension can compensate for lower performance in another. Other weighting and aggregation approaches, in turn, may produce different comparisons among municipalities, even when the indicators and underlying data remain the same.

This methodological dependence makes it necessary to examine the stability of municipal positions, since different choices of weights and aggregation methods may yield different rankings from the same indicators. The first research question is therefore: to what extent do these choices alter the ranking of municipalities? Furthermore, comparing the rankings with information not used in their construction makes it possible to investigate how they relate to a health outcome associated with sanitation conditions. The second question is thus: does any of the rankings show a stronger association with the hospitalization rate for diseases related to inadequate environmental sanitation (DRSAI) than the ABES ranking?

Accordingly, this study compares 16 weighting and aggregation configurations applied to the five sanitation indicators for the 505 municipalities in Northeast Brazil included in the 2026 ABES Ranking. First, it examines agreement among the rankings and changes in municipal positions. Next, it compares the association of each configuration with the DRSAI hospitalization rate, using the ABES ranking as a reference. Finally, an exploratory analysis identifies the municipalities that remain in the worst-performing quartile across all nine rankings retained after the comparison of their associations with health. The study thus investigates both the sensitivity of municipal positions to methodological choices and the existence of shared priorities across different rankings.